

\documentclass[11pt,a4paper,roman]{moderncv}      
\usepackage[italian]{babel}



\moderncvstyle{classic}                            
\moderncvcolor{green}                            

% character encoding
\usepackage[utf8]{inputenc}                     

% adjust the page margins
\usepackage[scale=0.75]{geometry}

% personal data
\name{Alessio}{Rovere}
\name{GeoMarine Engineering s.r.l.}{}
\email{alessio.rovere@unive.it} 
\address{Via delle Lagune 45 – 30172 Venezia (VE)}


\begin{document}

\recipient{A:}{Gruppo di ricerca Ambientale \\ Corso di Geografia Fisica e Geomorfologia}
\date{\today}
\makelettertitle

\noindent
\textbf{Oggetto:} \textit{Incarico per la valutazione della massima ingressione del runup costiero nel litorale di Borghetto Santo Spirito (SV)} \\[0.5em]

\noindent
\textbf{Venezia, \today} \\[1em]

Alla cortese attenzione del gruppo di consulenza ambientale, \\[0.5em]

La società \textbf{GeoMarine Engineering s.r.l.}, nell’ambito delle proprie attività di supporto tecnico al \textbf{Comune di Borghetto Santo Spirito (SV)}, conferisce al gruppo di consulenza ambientale l’incarico di elaborare una \textbf{valutazione della massima ingressione del runup costiero} durante un evento di mareggiata estrema, analogo a quello verificatosi nell’autunno 2018.

\section*{1. Scopo dell’incarico}
Il presente incarico si inserisce nell’ambito delle attività promosse dal \textbf{Comune di Borghetto Santo Spirito (SV)} per il miglioramento della \textbf{resilienza costiera a eventi meteoromarini estremi} e per la predisposizione di \textbf{mappe di rischio e vulnerabilità} da utilizzare come supporto alla pianificazione territoriale e alla gestione delle emergenze.

L’obiettivo dell’attività è la \textbf{valutazione quantitativa della massima ingressione del runup costiero} lungo la fascia litoranea comunale, mediante l’applicazione di un \textbf{ensemble di modelli empirici di runup} su dati di riferimento forniti.  
L’analisi dovrà consentire di individuare le aree potenzialmente soggette a \textbf{inondazione costiera durante eventi di mareggiata estrema}, analoghi a quello verificatosi nell’autunno 2018.

I risultati dovranno essere presentati in forma \textbf{cartografica e sintetica}, corredati da una discussione critica dei valori stimati e delle principali incertezze del metodo.  
Le informazioni prodotte potranno contribuire alla definizione di strumenti di pianificazione costiera, alla valutazione del rischio per le infrastrutture e al supporto delle strategie di adattamento locale ai cambiamenti climatici.

\section*{2. Dati e materiale di riferimento}
Per lo svolgimento dell’incarico saranno resi disponibili tutti i materiali necessari all’elaborazione e alla rappresentazione dei risultati, attraverso una \textbf{repository dedicata sulla pagina Moodle del corso}.

Il materiale comprende:
\begin{itemize}
  \item un \textbf{Digital Elevation Model (DEM)} dell’area di studio, con risoluzione adeguata alla scala dell’analisi costiera;
  \item una \textbf{ortofoto ad alta risoluzione}, utile per la rappresentazione e l’interpretazione delle forme costiere;
  \item un \textbf{notebook Jupyter} contenente i modelli empirici di runup, le funzioni di calcolo e gli esempi di visualizzazione dei risultati.
\end{itemize}

Tutti i file forniti dovranno essere utilizzati esclusivamente come \textbf{base di lavoro per l’elaborazione e la mappatura del runup}, senza ulteriori modifiche ai dataset originali.

Per agevolare l’utilizzo del materiale e lo sviluppo dell’elaborato, sarà resa disponibile una \textbf{playlist di tutorial video su YouTube} con istruzioni operative passo–passo, disponibile sulla pagina Moodle del corso.

\section*{3. Attività richieste}
Il gruppo di consulenza è incaricato di svolgere le attività elencate di seguito, seguendo le istruzioni contenute nel materiale didattico fornito e nella documentazione tecnica allegata alla repository del progetto.

\begin{enumerate}
  \item \textbf{Calcolo del runup costiero.}  
  Utilizzare il \textbf{notebook Jupyter fornito} per calcolare i valori di \textit{runup massimo} attraverso l’applicazione di un \textbf{ensemble di modelli empirici} disponibili nel codice.  
  I risultati numerici dovranno essere verificati e, se necessario, confrontati con le formule di riferimento citate nel notebook.

  \item \textbf{Elaborazione e rappresentazione cartografica.}  
  Utilizzare i valori di runup ottenuti per generare, tramite \textbf{QGIS} o software equivalente, mappe delle aree potenzialmente soggette a \textbf{inondazione costiera}.  
  Le mappature dovranno essere sovrapposte al DEM e all’ortofoto forniti, includendo elementi di riferimento (scala, nord, legenda, coordinate) e una chiara indicazione delle soglie di runup considerate.

  \item \textbf{Redazione della relazione tecnica.}  
  Preparare una relazione sintetica (massimo 5 pagine, bibliografia inclusa) secondo lo schema seguente:
  \begin{itemize}
    \item \textbf{Introduzione} – contesto generale, obiettivi, e riferimenti bibliografici scientifico-tecnici sulle dinamiche costiere dell’area ligure e di Borghetto S.S.;
    \item \textbf{Area di studio} – descrizione del contesto geografico e geomorfologico, con una mappa realizzata in QGIS;
    \item \textbf{Metodi e dati} – spiegazione delle formule di runup utilizzate e dei dati di input;
    \item \textbf{Risultati e discussione} – presentazione dei risultati (anche attraverso mappature ad hoc) e valutazione critica delle incertezze;
    \item \textbf{Conclusioni e raccomandazioni} – considerazioni sul rischio costiero e possibili misure di mitigazione.
  \end{itemize}
  La relazione dovrà essere redatta in stile tecnico–scientifico e corredata da citazioni in testo e da una bibliografia finale coerente con gli standard accademici.  

  Una \textbf{guida alla redazione della relazione tecnico–scientifica} è disponibile sulla pagina Moodle del corso.
\end{enumerate}

\section*{4. Consegna e tempistiche}
Lo svolgimento dell’attività e la consegna della relazione tecnica seguiranno la seguente pianificazione:

\begin{itemize}
  \item \textbf{12 e 13 gennaio 2026 – Sessioni di lavoro.}  
  Durante queste due giornate (6 ore ciascuna) i gruppi lavoreranno sul proprio elaborato, potendo interagire direttamente con il committente per domande tecniche, chiarimenti metodologici o revisione parziale dei risultati.  
  \textbf{Le richieste di feedback o aiuto al committente potranno essere formulate esclusivamente durante queste sessioni.}  
  Eventuali richieste via e-mail non verranno prese in considerazione, al fine di garantire equità tra tutti i gruppi di lavoro.

  \item \textbf{13 gennaio 2026 – Consegna della relazione preliminare (entro le ore 23:59).}  
  Ogni gruppo dovrà caricare su \textbf{piattaforma Moodle} una prima versione completa della relazione tecnica in formato PDF.  
  \textbf{Non sono ammesse consegne via e-mail o attraverso altri canali.}  
  La versione preliminare sarà oggetto di discussione e revisione nella fase successiva.

  \item \textbf{14 gennaio 2026 – Revisione e discussione guidata.}  
  Nel corso della giornata, i gruppi presenteranno sinteticamente il proprio elaborato e riceveranno indicazioni mirate per il miglioramento della relazione, con particolare attenzione alla coerenza metodologica, all’uso dei modelli di runup e alla qualità della rappresentazione cartografica.

  \item \textbf{15 gennaio 2026 – Consegna finale (entro le ore 23:59).}  
  La versione definitiva della relazione dovrà essere caricata obbligatoriamente su \textbf{Moodle} entro questo termine, come \textbf{file PDF unico per gruppo} (massimo 3 componenti).  
  Dopo tale scadenza, \textbf{ogni giorno di ritardo comporterà la sottrazione di 1 punto} dal punteggio finale assegnato all’elaborato.
\end{itemize}

Si ricorda che \textbf{non saranno accettate consegne via e-mail} né richieste di assistenza successive alle sessioni di lavoro in classe (12–14 gennaio 2026).  
Eventuali dubbi o necessità di chiarimento dovranno essere affrontati esclusivamente durante le giornate di lavoro supervisionato.  
La consegna oltre i termini stabiliti comporterà automaticamente l’applicazione delle penalità previste.  

\section*{5. Proprietà intellettuale e originalità}
La relazione tecnica dovrà essere \textbf{interamente originale} e redatta autonomamente dal gruppo incaricato.  
Ogni utilizzo di testi, dati, figure o mappe provenienti da altre fonti dovrà essere accompagnato da una \textbf{citazione esplicita in testo} e da un corretto riferimento bibliografico nel formato indicato nel \textit{Manuale di scrittura tecnico–scientifica} disponibile nella pagina Moodle del corso.

Tutti i contenuti dovranno rispettare le buone pratiche di redazione tecnica, che comprendono:
\begin{itemize}
  \item l’utilizzo di un linguaggio \textbf{chiaro, sintetico e privo di ambiguità};
  \item la \textbf{strutturazione logica} del testo in sezioni coerenti con il metodo scientifico (introduzione, materiali e metodi, risultati, discussione, conclusioni);
  \item la \textbf{presentazione accurata delle figure e delle tabelle}, con didascalie complete e riferimenti nel testo;
  \item la corretta indicazione delle \textbf{unità di misura, scale e coordinate} nelle rappresentazioni cartografiche;
  \item la \textbf{tracciabilità delle fonti} di dati, immagini o citazioni, garantendo la trasparenza e la verificabilità del lavoro svolto.
\end{itemize}

Sono espressamente vietati il \textbf{plagio}, la riproduzione non autorizzata di materiale di terzi o il riutilizzo di elaborati provenienti da altri gruppi o anni accademici.  
La violazione di tali norme comporterà la \textbf{nullità dell’incarico e la valutazione automatica di insufficienza}.

Si raccomanda di utilizzare gli strumenti indicati nel manuale per il controllo della coerenza testuale e delle citazioni bibliografiche, e di verificare la qualità formale dell’elaborato prima della consegna.

\section*{6. Divieto di subappalto e responsabilità}
Il presente incarico è conferito in via esclusiva al gruppo di consulenza assegnato e \textbf{non può essere subappaltato, delegato o condiviso} con soggetti esterni al gruppo.  
Ogni componente del gruppo è responsabile in solido della correttezza, completezza e originalità del lavoro svolto.

Non è consentita la collaborazione con altri gruppi o con persone non iscritte al corso, né lo scambio di codici, dati, elaborati o grafici prodotti nell’ambito dell’attività.  
Tali restrizioni mirano a garantire \textbf{pari condizioni di lavoro e valutazione} per tutte le studentesse e tutti gli studenti coinvolti.

Eventuali collaborazioni esterne — ad esempio, per l’uso di dati o riferimenti aggiuntivi provenienti da enti o progetti scientifici — dovranno essere \textbf{esplicitamente dichiarate e approvate} dal docente responsabile prima della consegna della relazione finale.

I materiali forniti nella repository ufficiale (DEM, ortofoto, notebook e tutorial) sono di uso strettamente didattico e non possono essere utilizzati o ridistribuiti per fini diversi da quelli previsti dall’incarico.  
Il riutilizzo di tali materiali in altri contesti richiede il consenso scritto del docente o del titolare dei diritti.

La violazione di una o più di queste condizioni potrà comportare la \textbf{revoca dell’incarico} e l’attribuzione di una valutazione negativa, indipendentemente dalla qualità tecnica dell’elaborato consegnato.

\section*{7. Criteri di valutazione tecnica}
La valutazione del lavoro sarà effettuata in base a criteri tecnico–scientifici e formali, volti a premiare la qualità metodologica, la chiarezza espositiva e la coerenza complessiva dell’elaborato.  
Il punteggio complessivo massimo attribuibile è di \textbf{30 punti}, che concorrerà alla determinazione del voto finale dell’esercitazione.

I principali criteri di valutazione sono i seguenti:

\begin{itemize}
  \item \textbf{Contenuto scientifico e metodologico} – correttezza nell’applicazione dei modelli di runup, qualità dei dati utilizzati e coerenza dei risultati con gli obiettivi dell’incarico;
  \item \textbf{Introduzione e inquadramento bibliografico} – capacità di contestualizzare il lavoro rispetto alla letteratura scientifica e ai rapporti tecnici disponibili, includendo riferimenti pertinenti e aggiornati;
  \item \textbf{Chiarezza e struttura della relazione} – organizzazione logica del testo, uso di un linguaggio tecnico appropriato, coerenza tra sezioni e fluidità nella presentazione dei risultati;
  \item \textbf{Qualità grafica e cartografica} – accuratezza delle figure e delle mappe prodotte, presenza di legenda, scala, coordinate e didascalie esplicative; leggibilità e pertinenza delle rappresentazioni;
  \item \textbf{Capacità di analisi e discussione critica} – interpretazione consapevole dei risultati, riconoscimento dei limiti metodologici e delle incertezze, formulazione di raccomandazioni operative;
  \item \textbf{Originalità e correttezza delle fonti} – rispetto delle buone pratiche di citazione, assenza di plagio e uso corretto dei dati e materiali forniti;
  \item \textbf{Rispetto delle tempistiche e delle istruzioni operative} – puntualità nelle consegne, conformità al formato richiesto e osservanza delle regole del presente incarico.
\end{itemize}

Un elaborato potrà essere considerato \textbf{insufficiente} qualora presenti gravi errori metodologici, carenze nella bibliografia o nella rappresentazione dei risultati, oppure in caso di plagio o mancata attribuzione delle fonti.

\section*{8. Corrispettivo, condizioni e valore formativo}
Il presente incarico rientra tra le attività previste dal corso universitario \textit{Geografia Fisica e Geomorfologia} del Corso di Laurea in Scienze Ambientali dell’\textbf{Università Ca’ Foscari Venezia}, a.a. 2025/2026.  
Si tratta di una \textbf{simulazione didattica} che riproduce un incarico tecnico reale, finalizzata allo sviluppo di competenze operative, analitiche e comunicative nel campo della geomorfologia costiera.

Non è previsto alcun \textbf{corrispettivo economico}, trattandosi di attività formativa obbligatoria. Tuttavia, la qualità dell’elaborato e il rispetto delle specifiche tecniche e metodologiche incideranno direttamente sulla \textbf{valutazione complessiva dell’esercitazione}, che contribuirà per \textbf{un terzo (1/3) al voto finale del corso}.

L’attività mira a fornire agli studenti e alle studentesse l’opportunità di:
\begin{itemize}
  \item applicare conoscenze teoriche alla \textbf{risoluzione di un problema tecnico reale}, con riferimento alla gestione del rischio costiero;
  \item sperimentare strumenti digitali di analisi e modellazione (Jupyter Notebook, QGIS, GitHub);
  \item produrre un elaborato tecnico conforme agli \textbf{standard professionali di ingegneria costiera e ambientale};
  \item sviluppare capacità di \textbf{lavoro di gruppo, organizzazione e comunicazione scientifica}.
\end{itemize}

Il lavoro svolto e i risultati conseguiti resteranno di esclusiva proprietà didattica del corso e potranno essere utilizzati a fini di documentazione o presentazione accademica, nel rispetto della normativa vigente in materia di diritti d’autore.

Si invita ogni gruppo a considerare l’attività non solo come un esercizio accademico, ma come una \textbf{simulazione realistica di consulenza tecnica ambientale}, in cui accuratezza, chiarezza e rispetto delle procedure costituiscono parte integrante della valutazione finale.  

\vspace{1em}
Rimaniamo a disposizione per eventuali chiarimenti tecnici o metodologici durante lo svolgimento dell’attività.

\vspace{1em}
Cordiali saluti, \\[1em]
\textbf{Dott.. Alessio Rovere} \\
\textit{Responsabile tecnico – GeoMarine Engineering s.r.l.}
\end{document}
